\subsection{测度代数}

设 G 是局部紧群。以后始终约定，G 上的测度 \(\mu_{1},\ldots,\mu_{n}\) 称为可卷积的，是指它们对于映射

\[
(x _ {1}, x _ {2}, \ldots , x _ {n}) \mapsto x _ {1} x _ {2} \cdot \cdot \cdot x _ {n};
\]

卷积积 \(\ast_{i}\mu_{i}\) 总是通过这个映射定义。若 \(s\in G, t\in G\)，则

\[
\varepsilon_ {s} * \varepsilon_ {t} = \varepsilon_ {s t}.\tag{1}
\]

若 \(s \in \mathbf{G}\) 且 \(\mu \in \mathcal{M}(\mathbf{G})\)，则

\begin{enumerate}
\def\labelenumi{(\arabic{enumi})}
\setcounter{enumi}{1}
\tightlist
\item
\end{enumerate}

\[
\begin{array}{l} \varepsilon_ {s} * \mu = \gamma (s) \mu \\ \mu * \varepsilon_ {s} = \delta (s ^ {- 1}) \mu \end{array}\tag{3}
\]

根据 §1 No.~1 的例 3。若 G 是交换群，则 \(\mu_{1}\) 与 \(\mu_{2}\) 可卷积等价于 \(\mu_{2}\) 与 \(\mu_{1}\) 可卷积，并且此时有 \(\mu_{1} * \mu_{2} = \mu_{2} * \mu_{1}\)。当 G 非交换时，可能出现 \(\mu_{1}\) 与 \(\mu_{2}\) 可卷积，而 \(\mu_{2}\) 与 \(\mu_{1}\) 不可卷积的情形（习题 12）。

命题 1. --- 设 G 是局部紧群，\(\lambda, \mu, \nu\) 是 G 上均 \(\neq 0\) 的测度。

\begin{enumerate}
\def\labelenumi{(\roman{enumi})}
\tightlist
\item
  若 \(\lambda, \mu, \nu\) 可卷积，则 \(\lambda\) 与 \(\mu\)、\(|\lambda| * |\mu|\) 与 \(\nu\)、\(\mu\) 与 \(\nu\)、\(\lambda\) 与 \(|\mu| * |\nu|\) 也都可卷积，并且有
\end{enumerate}

\[
\lambda * \mu * \nu = (\lambda * \mu) * \nu = \lambda * (\mu * \nu).
\]

\begin{enumerate}
\def\labelenumi{(\roman{enumi})}
\setcounter{enumi}{1}
\tightlist
\item
  若 \(\lambda\) 与 \(\mu\) 可卷积，且 \(|\lambda| * |\mu|\) 与 \(\nu\) 也可卷积，则 \(\lambda, \mu, \nu\) 可卷积。若 \(\mu\) 与 \(\nu\) 可卷积，且 \(\lambda\) 与 \(|\mu| * |\nu|\) 也可卷积，同样如此。
\end{enumerate}

这由 §1 No.~2 的命题 1 得出。

\pandocbounded{\includegraphics[keepaspectratio]{images/part001_0fdec0aa.jpg}}

测度 \(\lambda, \mu, \nu\) 定义在 \(G\) 上可能存在，使卷积积 \(\lambda * \mu\)、\((\lambda * \mu) * \nu\)、\(\mu * \nu\)、\(\lambda * (\mu * \nu)\) 都有定义，但仍有 \((\lambda * \mu) * \nu \neq \lambda * (\mu * \nu)\)（参见习题 4）。

设 \(\rho\) 是取有限值且 \(>0\) 的下半连续函数，定义在 \(\mathbf{G}\) 上，满足 \(\rho(st) \leqslant \rho(s)\rho(t)\) 对所有 \(s, t\) 在 \(\mathbf{G}\) 中成立。记 \(\mathcal{M}^{\rho}(\mathbf{G})\) 为由测度 \(\lambda\) 构成的向量空间，其中这些测度定义在 \(\mathbf{G}\) 上，且 \(\rho\) 对 \(\lambda\) 可积，并记 \(\|\lambda\|_{\rho}\)（或简写为 \(\|\lambda\|\)）为该空间上的范数 \(\int_{\mathbf{G}} \rho(s)d|\lambda|(s)\)。当 \(\rho = 1\) 时，就得到 \(\mathcal{M}^{1}(\mathbf{G})\)，即 \(\mathbf{G}\) 上有界测度的集合。

命题 2. --- (i) \(\mathcal{M}^{\rho}(\mathrm{G})\) 中任意两个元素都可卷积。

\begin{enumerate}
\def\labelenumi{(\roman{enumi})}
\setcounter{enumi}{1}
\item
  对于卷积和范数 \(\| \lambda \|\)，\(\mathcal{M}^{\rho}(\mathbf{G})\) 是完备赋范代数，并以 \(\varepsilon_{e}\) 为单位元。
\item
  \(\mathcal{C}'(\mathbf{G})\) 是 \(\mathcal{M}^{\rho}(\mathbf{G})\) 的子代数。
\end{enumerate}

设 \(\lambda, \mu\) 属于 \(\mathcal{M}^{\rho}(\mathbf{G})\)，我们来证明 \(\lambda\) 与 \(\mu\) 可卷积。设 \(f \in \mathcal{K}_{+}(\mathbf{G})\)。由于 \(\rho\) 是 \(>0\) 且下半连续的，存在常数 \(k > 0\) 使得 \(f \leqslant k\rho\)。于是

\[
\begin{array}{r l} & {\int^ {*} f (s t) d | \lambda | (s) d | \mu | (t) \leqslant k \int^ {*} \rho (s t) d | \lambda | (s) d | \mu | (t)} \\ & {\qquad \leqslant k \int^ {*} \rho (s) \rho (t) d | \lambda | (s) d | \mu | (t)} \\ & {\qquad = k \bigg (\int^ {*} \rho (s) d | \lambda | (s) \bigg) \bigg (\int^ {*} \rho (t) d | \mu | (t) \bigg)} \end{array}
\]

（第 V 章 §8 No.~3 命题 8 的推论 1）。因此，\((s,t) \mapsto f(st)\) 是 \((\lambda \otimes \mu)\)-可积的，从而 \(\lambda\) 与 \(\mu\) 可卷积。另一方面，利用第 V 章（§1 命题 4、§6 命题 2、§8 命题 8 的推论 1）以及 \((s,t) \mapsto \rho(s)\rho(t)\) 在 \(G \times G\) 上下半连续这一事实，有

\[
\begin{array}{r l} & {\int_ {\mathrm{G}} ^ {*} \rho (s) d | \lambda * \mu | (s) = \int_ {\mathrm{G}} ^ {\bullet} \rho (s) d | \lambda * \mu | (s)} \\ & {\quad \leqslant \int_ {\mathrm{G} \times \mathrm{G}} ^ {\bullet} \rho (s t) d | \lambda | (s) d | \mu | (t) \leqslant \int_ {\mathrm{G} \times \mathrm{G}} ^ {\bullet} \rho (s) \rho (t) d | \lambda | (s) d | \mu | (t)} \\ & {\quad = \int_ {\mathrm{G} \times \mathrm{G}} ^ {*} \rho (s) \rho (t) d | \lambda | (s) d | \mu | (t) = \| \lambda \| \cdot \| \mu \|.} \end{array}
\]

可见 \(\lambda * \mu \in \mathcal{M}^{\rho}(\mathrm{G})\)，且 \(\| \lambda * \mu \| \leqslant \| \lambda \| \cdot \| \mu \|\)。根据命题 1，\(\mathcal{M}^{\rho}(\mathrm{G})\) 是一个代数。映射 \(\lambda \mapsto \rho \cdot \lambda\) 是等距线性映射 \(\theta\)，将 \(\mathcal{M}^{\rho}(\mathrm{G})\) 映入 \(\mathcal{M}^{1}(\mathrm{G})\)；若 \(\mu \in \mathcal{M}^{1}(\mathrm{G})\)，则 \(1 / \rho\) 局部有界且上半连续，局部 \(\mu\)-可积，而 \(\rho\) 对 \((1 / \rho) \cdot \mu\) 可积，故 \((1 / \rho) \cdot \mu \in \mathcal{M}^{\rho}(\mathrm{G})\)；这证明了 \(\theta\) 是满射；因此 \(\mathcal{M}^{\rho}(\mathrm{G})\) 是完备赋范代数。最后，显然 \(\varepsilon_{e}\) 是 \(\mathcal{M}^{\rho}(\mathrm{G})\) 的单位元，且 \(\mathcal{C}'(\mathrm{G})\) 是 \(\mathcal{M}^{\rho}(\mathrm{G})\) 的子代数（§1 No.~4 命题 5 的推论）。

若 \(\rho = 1\)，命题 2 的 (i) 和 (ii) 也由 §1 的命题 2 得出。

命题 3. --- 设 \(\mu_1, \ldots, \mu_n\) 是 G 上的测度。若这些 \(\mu_i\) 中至多一个没有紧支集，则 \(\mu_i\) 可卷积。

事实上，令 \(S_{i}\) 为 \(\mu_{i}\) 的支集，并设 \(S_{i}\) 对 \(i \neq i_{0}\) 是紧的。令 K 为 G 的一个紧子集。满足 \((x_{1}, \ldots, x_{n}) \in \prod_{i} S_{i}\) 且 \(x_{1}x_{2} \cdots x_{n} \in K\) 的集合是紧的，因为条件 \(x_{i} \in S_{i}\) 对所有的 \(i, x_{1}x_{2} \cdots x_{n} \in K\) 均成立，蕴含

\[
x _ {i _ {0}} \in \mathrm{S} _ {i _ {0} - 1} ^ {- 1} \dots \mathrm{S} _ {1} ^ {- 1} \mathrm{KS} _ {n} ^ {- 1} \dots \mathrm{S} _ {i _ {0} + 1} ^ {- 1}.
\]

因此 \(\mu_{i}\) 可卷积（§1 No.~4 命题 4）。

命题 4. --- 映射 \((\lambda, \mu) \mapsto \lambda * \mu\)（分别为 \((\lambda, \mu) \mapsto \mu * \lambda)\)），其中 \(\lambda \in \mathcal{C}'(G)\)、\(\mu \in \mathcal{M}(G)\)，在 \(\mathcal{M}(G)\) 上定义了关于代数 \(\mathcal{C}'(G)\) 的左（分别为右）模结构。

这由命题 1 和命题 3 得出。

命题 5. --- 设 \(\lambda\) 是 G 上的左（分别为右）Haar 测度，且 \(\mu \in \mathcal{M}^1(\mathrm{G})\)。则 \(\mu\) 与 \(\lambda\)（分别为 \(\lambda\) 与 \(\mu\)）可卷积，且 \(\mu * \lambda = \mu(1)\lambda\)（分别为 \(\lambda * \mu = \mu(1)\lambda\)）。

不妨设 \(\mu \geqslant 0\)。令 \(f \in \mathcal{K}_{+}(\mathbf{G})\)。当 \(\lambda\) 是左 Haar 测度时，

\[
\int^ {*} d \mu (x) \int^ {*} f (x y) d \lambda (y) = \int^ {*} d \mu (x) \int f (y) d \lambda (y) = \lambda (f) \| \mu \|,
\]

因此函数 \((x,y)\mapsto f(xy)\) 是 \((\mu \otimes \lambda)\)-可积的，并且它关于 \(\mu \otimes \lambda\) 的积分为 \(\lambda (f)\| \mu \|\)。当 \(\lambda\) 是右 Haar 测度时，论证类似。

命题 6. --- 设 \(\mu\) 与 \(\nu\) 是 \(\mathbf{G}\) 上的两个可卷积测度。设 \(\chi\) 是 \(\mathbf{G}\) 在 \(\mathbf{C}^*\) 中的连续表示。则 \(\chi \cdot \mu\) 与 \(\chi \cdot \nu\) 可卷积，且 \((\chi \cdot \mu) * (\chi \cdot \nu) = \chi \cdot (\mu * \nu)\)。

令 \(f \in \mathcal{K}(\mathbf{G})\)。则 \(f\chi \in \mathcal{K}(\mathbf{G})\)，从而函数

\[
(x, y) \mapsto f (x y) \chi (x y) = f (x y) \chi (x) \chi (y)
\]

在 \(\mathbf{G} \times \mathbf{G}\) 上关于 \(\mu \otimes \nu\) 可积；因此函数 \((x, y) \mapsto f(xy)\) 关于 \((\chi \cdot \mu) \otimes (\chi \cdot \nu)\) 可积；所以 \(\chi \cdot \mu\) 与 \(\chi \cdot \nu\) 可卷积。此外，

\[
\begin{array}{r l} \langle \chi \cdot \mu * \chi \cdot \nu , f \rangle & = \int f (x y) \chi (x) \chi (y) d \mu (x) d \nu (y) \\ & = \int (f \chi) (x y) d \mu (x) d \nu (y) = \langle \mu * \nu , \chi f \rangle , \end{array}
\]

由此得到 \((\chi \cdot \mu) * (\chi \cdot \nu) = \chi \cdot (\mu * \nu)\)。

命题 7. --- 设 G 和 G' 是两个局部紧群，u 是 G 在 G' 中的连续表示，\(\mu_{1},\ldots,\mu_{n}\) 是 G 上均 \(\neq0\) 的测度。则下列断言等价：

\begin{enumerate}
\def\labelenumi{(\roman{enumi})}
\item
  \(u\) 对 \(\mu_{i}\) 是适当的，对所有 \(i\) 均如此，且测度 \(u(|\mu_i|)\) 可卷积；
\item
  \(\mu_{i}\) 可卷积，且 \(u\) 对 \(*_{i}(|\mu_{i}|)\) 是适当的
\end{enumerate}

满足这些条件时，

\[
\underset {i} {\ast} u (\mu_ {i}) = u \Big (\underset {i} {\ast} \mu_ {i} \Big).
\]

这由 §1 No.~2 命题 1 的推论得出。

推论。 --- 设 G 是局部紧群，\(\mu_{1},\ldots,\mu_{n}\) 是 G 上的测度。序列 \((\mu_{i})_{1\leqslant i\leqslant n}\) 可卷积的充要条件是序列 \((\breve{\mu}_{n-i})_{0\leqslant i\leqslant n-1}\) 可卷积；在此情形下，

\[
\left(\mu_ {1} * \dots * \mu_ {n}\right) ^ {\vee} = \stackrel {\vee} {\mu} _ {n} * \dots * \stackrel {\vee} {\mu} _ {1}.
\]

这是将同构 \(x \mapsto x^{-1}\) 从 \(G\) 到对立群应用于命题 7 的结果。

